\section{Counterfactuals and policy options}\label{sec:policy}

Table~\ref{tab:counterfactuals} changes one element of the end-2025 balance sheet at a time and reports the combined metric. Unless stated, debt/GDP is held at its baseline, so that each scenario isolates a change in composition. In the no-QE scenario, the Fed holds only as many Treasuries as it has currency outstanding, as before 2008; the rest return to private holders pro rata to the SOMA maturity structure, and reserves and reverse repos go to zero. Ending interest on reserves, or tiering, stops all, half or a quarter of the \$2.85tn of reserves earning interest, which moves them from overnight interest-bearing debt to the zero-interest base that inflation erodes. In the bills scenario, the Fed sells coupons pro rata to private holders and buys bills from them until half its Treasuries are bills, a change of about \$1.9tn. In the bill-share scenarios, the Treasury raises bills to 25\% or 30\% of marketable debt and shrinks all other securities pro rata.

\begin{table}[htbp]
\centering
\caption{End-2025 counterfactuals}\label{tab:counterfactuals}
{\small +1pp permanent rate rise, H = 10; baseline $\hat\varphi$ = 0\par}\smallskip
\footnotesize
\begin{tabularx}{\textwidth}{>{\raggedright\arraybackslash}p{0.3\textwidth}>{\centering\arraybackslash}X>{\centering\arraybackslash}X>{\centering\arraybackslash}X>{\centering\arraybackslash}X>{\centering\arraybackslash}X}
\toprule
Scenario & $\varphi^*$ & Gap & Inflation to cover gap (pp/yr) & Change vs. baseline & One-time jump (\%) \\
\midrule
Baseline (consolidated; g = 3.8\%, SEP) & 0.46 & 0.46 & 1.00 & — & 3.26 \\
\emph{Fiscal policy} &  &  &  &  &  \\
Fiscal response at the literature's pre-2004 value ($\hat\varphi$ = 0.39) & 0.46 & 0.07 & 0.16 & $-$84\% & 0.51 \\
Fiscal response 0.25 (intermediate) & 0.46 & 0.21 & 0.46 & $-$54\% & 1.50 \\
\emph{Central-bank balance sheet} &  &  &  &  &  \\
No QE: Fed Treasuries = currency, no reserves & 0.46 & 0.46 & 0.83 & $-$18\% & 3.07 \\
QT: reserves to \$1.9tn, Treasuries back to private & 0.46 & 0.46 & 0.94 & $-$6\% & 3.20 \\
Ending interest on reserves & 0.46 & 0.46 & 0.67 & $-$33\% & 2.82 \\
Tiering: 50\% of reserves unremunerated & 0.46 & 0.46 & 0.82 & $-$19\% & 3.04 \\
Tiering: 25\% of reserves unremunerated & 0.46 & 0.46 & 0.90 & $-$10\% & 3.15 \\
Fed Treasury portfolio 50\% bills (from 5.5\%) & 0.46 & 0.46 & 0.89 & $-$12\% & 3.14 \\
\emph{Treasury debt management} &  &  &  &  &  \\
Bill share 25\% of marketable debt (from 21.6\%) & 0.46 & 0.46 & 1.06 & +6\% & 3.31 \\
Bill share 30\% & 0.46 & 0.46 & 1.15 & +15\% & 3.39 \\
Terms out: 10\% of bills $\to$ 10-year notes & 0.46 & 0.46 & 0.93 & $-$7\% & 3.18 \\
\emph{Markets and macro} &  &  &  &  &  \\
Convenience yield $-$50bp & 0.50 & 0.50 & 1.08 & +8\% & 3.51 \\
Convenience yield $-$100bp & 0.53 & 0.53 & 1.15 & +14\% & 3.73 \\
Trend growth 4\% & 0.43 & 0.43 & 0.95 & $-$5\% & 3.07 \\
Trend growth 3.5\% & 0.50 & 0.50 & 1.08 & +8\% & 3.55 \\
$\psi$ = 2bp (CBO) & 0.38 & 0.38 & 0.82 & $-$18\% & 2.67 \\
$\psi$ = 4.5bp & 0.55 & 0.55 & 1.20 & +20\% & 3.90 \\
\emph{Accounting benchmark} &  &  &  &  &  \\
Treasury's own liabilities (Fed ignored) & 0.46 & 0.46 & 1.02 & +1\% & 3.31 \\
\bottomrule
\end{tabularx}
\end{table}


Five results follow. First, the fiscal response is the dominant lever: a response at the literature's pre-2004 value (0.39) cuts the requirement by 84\%, and no balance-sheet policy comes close. Second, balance-sheet policy moves only the inflation layer. None of the central-bank or debt-management scenarios changes $\varphi^*$ when debt/GDP is held fixed; they change the price of the inflation route by 6--33\%.

Third, ending interest on reserves is the largest balance-sheet lever. It cuts the requirement by a third (36\% once debt/GDP falls with the interest-bearing stock), almost twice the effect of undoing QE, and tiering at 50\% or 25\% cuts it by 19\% or 10\%. In levels (Section~\ref{sec:level}), it lowers the inflation needed to hold debt/GDP stable from 4.2 to 2.3pp a year, because reserves stop costing interest (about \$104bn a year at the end-2025 rate of 3.65\%) and their growth becomes seigniorage. It is not a free lunch. Unremunerated reserves are a tax on banks, and in an ample-reserves framework the interest rate on reserves is also how the Fed sets the policy rate, so ending it would require far smaller reserves or binding reserve requirements. The ECB's move to stop remunerating minimum reserves in 2023 is a partial version of the same idea.

Fourth, a Fed shift to bills lengthens the consolidated structure. Moving half the Fed's Treasuries into bills cuts the requirement by 12\%: the Fed's own portfolio gets shorter, but the private sector ends up holding more coupons, and it is the private sector's holdings that matter. From the consolidated point of view, the Fed buying bills is the Treasury terming out. Fifth, more bills raise the price of the inflation route: a bill share of 25\% raises the requirement by 6\%, and 30\% by 15\%, while terming out a tenth of bills lowers it by 7\%.

When debt/GDP is allowed to change with the interest-bearing stock, the results are similar. No QE without the reserves that fund the Fed's MBS lowers $b$ from 0.95 to 0.90 and the requirement by 20\%; ending interest on reserves lowers $b$ to 0.86, $\varphi^*$ to 0.44 and the requirement by 36\%; tiering at 50\% lowers the requirement by 20\%. Losing the convenience yield, slower growth and a steeper debt-rate schedule are different in kind: they raise the threshold itself. Two limitations apply. The balance-sheet scenarios hold $r$ fixed and ignore the term-premium effects of QE, QT and changes in bill supply, and all scenarios are static comparisons.
